\documentclass[
  aps,
  pra,
  reprint,
  superscriptaddress,
  nofootinbib,
  longbibliography,
  floatfix,
  amsmath,
  amssymb
]{revtex4-2}

\usepackage[T1]{fontenc}
\usepackage[utf8]{inputenc}
\usepackage{lmodern}
\usepackage{amsthm}
\usepackage{bm}
\usepackage{graphicx}
\usepackage{booktabs}
\usepackage{array}
\allowdisplaybreaks

\theoremstyle{plain}
\newtheorem{theorem}{Theorem}
\newtheorem{lemma}[theorem]{Lemma}
\newtheorem{proposition}[theorem]{Proposition}
\newtheorem{corollary}[theorem]{Corollary}
\theoremstyle{definition}
\newtheorem{definition}[theorem]{Definition}
\theoremstyle{remark}
\newtheorem{remark}[theorem]{Remark}

\newcommand{\ket}[1]{\lvert #1\rangle}
\newcommand{\bra}[1]{\langle #1\rvert}
\newcommand{\braket}[2]{\langle #1\vert #2\rangle}
\newcommand{\Tr}{\operatorname{Tr}}
\newcommand{\SU}{\mathrm{SU}}
\newcommand{\SO}{\mathrm{SO}}
\newcommand{\UU}{\mathrm{U}}
\newcommand{\CP}{\mathbb{CP}}
\newcommand{\RR}{\mathbb{R}}
\newcommand{\CC}{\mathbb{C}}
\newcommand{\Id}{\mathbb{I}}
\newcommand{\Pauli}{\boldsymbol{\sigma}}

\raggedbottom
\begin{document}

\title{A Hopf-Bundle Holonomy Representation of Spin-Singlet Bell Correlations}

\author{Andrei Preece}
\email{andre.preece@web.de}
\affiliation{Independent Researcher, Berlin, Germany}
\author{Boris Batenin}
\email{batenin@web.de}
\affiliation{Independent Researcher, Berlin, Germany}
\date{April 2026}

\begin{abstract}
We give a conservative geometric representation of the spin-singlet Bell correlation
\(E(\bm a,\bm b)=-\bm a\cdot\bm b\) as the real part of a closed \(\UU(1)\) holonomy on the standard Hopf bundle \(S^1\to S^3\to S^2\) associated with normalized spinors and Bloch-sphere rays.  The construction uses the usual Berry connection and a closed loop on the Bloch sphere enclosing solid angle \(2(\pi-\theta)\), where \(\theta\) is the angle between the measurement settings.  It does not introduce local hidden variables, modify the Born probabilities, or change the no-signaling consequences of the singlet state.  We prove a restricted rigidity result: within this Hopf-bundle lune construction, an integer Chern charge reproduces the singlet kernel for all settings if and only if \(q=\pm1\), with the standard Hopf orientation written as \(q=+1\).  We also state the associated rotational symmetry, the reduced-state no-signaling property, and the standard CHSH/Tsirelson bound.  The result is a self-contained representation theorem linking the Pauli-algebra calculation, the Hopf fibration, and Berry holonomy without changing the operational content of quantum mechanics.
\end{abstract}

\maketitle

\section{Scope and Limitations}
\label{sec:purpose}

This paper is a quantum-foundations article about the geometry already present in the standard spin-\(1/2\) formalism.  The compact space used below is the unit spinor sphere \(S^3\subset\CC^2\) and its Hopf projection to the Bloch sphere \(S^2\); no assumption about cosmological spatial topology is used.

The central statement is the following restricted one:
\begin{quote}
The standard singlet correlation \(E(\bm a,\bm b)=-\cos\theta\), where \(\theta=\arccos(\bm a\cdot\bm b)\), admits a gauge-invariant closed-loop holonomy representation on the Hopf bundle.  The representation uses only the standard spinor geometry behind ordinary quantum mechanics.
\end{quote}

We explicitly do \emph{not} claim that this construction:
\begin{itemize}
\item provides a local hidden-variable model;
\item restores Bell-local factorization;
\item explains Bell violation by superluminal communication;
\item changes quantum mechanics, the Born rule, or the Hilbert-space tensor product;
\item derives a new experimental prediction distinct from the singlet-state prediction.
\end{itemize}

The contribution is structural.  It gives a clean geometric dictionary between three standard facts: (i) entangled spin correlations, (ii) the Hopf bundle and Berry connection of spinors, and (iii) the phase/holonomy character of the correlation kernel.  The rigidity claim is correspondingly internal to this dictionary; it is not a uniqueness theorem over all possible geometric encodings of the singlet kernel.

\section{Standard Algebraic Starting Point}
\label{sec:standard}

Let \(\bm a,\bm b\in S^2\) be two unit measurement directions and let
\begin{equation}
A_{\bm a}=\bm a\cdot\Pauli,\qquad B_{\bm b}=\bm b\cdot\Pauli,
\end{equation}
where \(\Pauli=(\sigma_x,\sigma_y,\sigma_z)\) are the Pauli matrices.  The spin singlet is
\begin{equation}
\ket{\Psi^-}=\frac{1}{\sqrt2}\bigl(\ket{01}-\ket{10}\bigr).
\end{equation}
The Bell correlation is
\begin{equation}
E(\bm a,\bm b)=
\bra{\Psi^-}(A_{\bm a}\otimes B_{\bm b})\ket{\Psi^-}.
\label{eq:Edef}
\end{equation}

\begin{proposition}[Singlet correlation]
For the singlet state,
\begin{equation}
E(\bm a,\bm b)=-\bm a\cdot\bm b=-\cos\theta,
\label{eq:singletcorr}
\end{equation}
where \(\theta\) is the angle between \(\bm a\) and \(\bm b\).
\end{proposition}

\begin{proof}
The singlet density matrix is
\begin{equation}
\rho_{\Psi^-}=\ket{\Psi^-}\bra{\Psi^-}
=\frac14\left(\Id\otimes\Id-
\sum_{i=1}^{3}\sigma_i\otimes\sigma_i\right).
\label{eq:singletdensity}
\end{equation}
Using \(\Tr(\sigma_i\sigma_j)=2\delta_{ij}\) and \(\Tr\sigma_i=0\),
\begin{align}
E(\bm a,\bm b)
&=\Tr\!\left[\rho_{\Psi^-}(\bm a\cdot\Pauli\otimes\bm b\cdot\Pauli)\right] \\
&=-\frac14\sum_{i,j,k} a_j b_k
\Tr(\sigma_i\sigma_j)\Tr(\sigma_i\sigma_k) \\
&=-\sum_i a_i b_i=-\bm a\cdot\bm b .
\end{align}
\end{proof}

This is the full physical content of the usual calculation.  The later holonomy construction must reproduce Eq.~\eqref{eq:singletcorr}, not replace it.

\section{Topology and Geometry of Spinors}
\label{sec:hopf}

A normalized spinor \(\psi\in\CC^2\) satisfies
\begin{equation}
\lVert\psi\rVert^2=|z_1|^2+|z_2|^2=1,
\end{equation}
so the space of normalized spinors is \(S^3\subset\CC^2\simeq\RR^4\).  Physical rays identify spinors differing by a phase, \(\psi\sim e^{i\alpha}\psi\), giving
\begin{equation}
\CP^1\simeq S^2.
\end{equation}
Thus spin-\(1/2\) pure states form the Hopf fibration
\begin{equation}
S^1\longrightarrow S^3\longrightarrow S^2.
\label{eq:hopf}
\end{equation}
The base \(S^2\) is the Bloch sphere.  The fiber \(S^1\) is the quantum phase.

\begin{remark}[Topology]
The relevant compact topology in this paper is not a cosmological \(S^3\).  It is the normalized spinor manifold \(S^3\subset\RR^4\).  It is compact, connected, simply connected, and boundaryless.  Its Hopf fiber has topology \(S^1\), with \(\pi_1(S^1)=\mathbb Z\), while \(\pi_1(S^3)=0\).  The base \(S^2\) is compact, connected, simply connected, and has \(\pi_1(S^2)=0\).
\end{remark}

A standard section of the Hopf bundle over the northern chart is
\begin{align}
\ket{\bm n}&=\begin{pmatrix}
\cos(\vartheta/2)\\
e^{i\varphi}\sin(\vartheta/2)
\end{pmatrix},\nonumber\\
\bm n&=(\sin\vartheta\cos\varphi,\sin\vartheta\sin\varphi,\cos\vartheta).
\label{eq:spinorsection}
\end{align}
The Berry connection on this chart is
\begin{equation}
A=-i\braket{\bm n}{d\bm n}
=\frac{1-\cos\vartheta}{2}\,d\varphi,
\label{eq:berryA}
\end{equation}
with curvature
\begin{equation}
F=dA=\frac12\sin\vartheta\,d\vartheta\wedge d\varphi.
\label{eq:berryF}
\end{equation}
For a closed loop \(\Gamma\subset S^2\), the holonomy is
\begin{equation}
U(\Gamma)=\exp\left(i\oint_\Gamma A\right)
=\exp\left(\frac{i}{2}\Omega[\Gamma]\right),
\label{eq:berryholonomy}
\end{equation}
where \(\Omega[\Gamma]\) is the oriented solid angle enclosed by \(\Gamma\), modulo \(4\pi\).

\section{Holonomy Representation of the Bell Kernel}
\label{sec:holonomy}

The singlet correlation is a real scalar function of the relative angle \(\theta\):
\begin{equation}
E(\theta)=-\cos\theta.
\end{equation}
A closed-loop phase can encode this function because
\begin{equation}
-\cos\theta=\operatorname{Re}\,e^{i(\pi-\theta)}.
\label{eq:phasekernel}
\end{equation}
The geometric task is therefore to associate to the pair \((\bm a,\bm b)\) a closed loop \(\Gamma_{ab}\) on the Bloch sphere with holonomy phase
\begin{equation}
\gamma_{ab}=\oint_{\Gamma_{ab}} A=\pi-\theta.
\label{eq:gammabell}
\end{equation}
Because \(\gamma=\Omega/2\), this is equivalent to choosing a loop with oriented solid angle
\begin{equation}
\Omega[\Gamma_{ab}]=2(\pi-\theta).
\label{eq:solidangle}
\end{equation}

\begin{definition}[Bell lune]
Given two points \(\bm a,\bm b\in S^2\) separated by angle \(\theta\), define a Bell lune \(\Gamma_{ab}\) as the oriented solid-angle class of closed curves satisfying
\begin{equation}
\Omega[\Gamma_{ab}]=2(\pi-\theta)\pmod{4\pi}.
\end{equation}
A canonical representative is the boundary of a spherical lune complementary to the minor geodesic separation of \(\bm a\) and \(\bm b\).  The curve is not an additional physical trajectory; by Stokes' theorem, only the oriented solid angle modulo \(4\pi\) enters the \(\UU(1)\) holonomy.
\end{definition}

\begin{theorem}[Holonomy representation]
For every pair of measurement directions \(\bm a,\bm b\in S^2\), the singlet Bell correlation can be written as
\begin{equation}
E(\bm a,\bm b)=\operatorname{Re}\,U(\Gamma_{ab})
=\operatorname{Re}\exp\left(i\oint_{\Gamma_{ab}}A\right)
=-\cos\theta,
\label{eq:holonomycorrelation}
\end{equation}
provided \(\Gamma_{ab}\) is chosen with solid angle \eqref{eq:solidangle}.
\end{theorem}

\begin{proof}
By Eq.~\eqref{eq:berryholonomy}, the closed-loop holonomy is \(U(\Gamma)=\exp(i\Omega[\Gamma]/2)\).  For the Bell lune, \(\Omega[\Gamma_{ab}]=2(\pi-\theta)\).  Therefore
\begin{equation}
U(\Gamma_{ab})=e^{i(\pi-\theta)}.
\end{equation}
Taking the real part gives \(\operatorname{Re}U=-\cos\theta\), which is exactly Eq.~\eqref{eq:singletcorr}.
\end{proof}


\subsection{Chern-charge rigidity of the holonomy construction}
\label{subsec:chernrigidity}

The preceding theorem uses the standard Hopf line bundle, whose curvature has unit Chern number.  It is useful to record exactly how much rigidity this assumption carries.  The following statement is internal to the Hopf-bundle lune construction; it is not a derivation of the Born rule from topology alone and does not exclude other, inequivalent geometric encodings.

\begin{theorem}[Chern-charge rigidity]
\label{thm:chernrigidity}
Let \(F\) be the Berry curvature of the standard Hopf line bundle, normalized by
\begin{equation}
\frac{1}{2\pi}\int_{S^2}F=1.
\end{equation}
For an integer Chern charge \(q\in\mathbb Z\), set
\begin{equation}
F_q=qF.
\end{equation}
For a Bell lune \(\Gamma_{ab}\) with
\begin{equation}
\Omega[\Gamma_{ab}]=2(\pi-\theta),
\qquad
\theta=\arccos(\bm a\cdot\bm b),
\end{equation}
define the corresponding real holonomy kernel by
\begin{equation}
E_q(\theta)=\operatorname{Re}\exp\left(iq\frac{\Omega[\Gamma_{ab}]}{2}\right).
\end{equation}
Then
\begin{equation}
E_q(\theta)=\cos\bigl(q(\pi-\theta)\bigr).
\end{equation}
The exact singlet condition
\begin{equation}
E_q(\theta)=-\cos\theta
\qquad
\text{for all }\theta\in[0,\pi]
\end{equation}
holds if and only if
\begin{equation}
q=\pm1.
\end{equation}
After fixing the standard orientation of the Hopf bundle, this is written as \(q=+1\).
\end{theorem}

\begin{proof}
For \(F_q=qF\), the closed-loop holonomy is
\begin{equation}
U_q(\Gamma_{ab})=
\exp\left(iq\frac{\Omega[\Gamma_{ab}]}{2}\right)
=\exp\bigl(iq(\pi-\theta)\bigr).
\end{equation}
Taking the real part gives
\begin{equation}
E_q(\theta)=\operatorname{Re}\,e^{iq(\pi-\theta)}
=\cos\bigl(q(\pi-\theta)\bigr).
\end{equation}

\textbf{Step 1: necessity of odd \(q\).}
At \(\theta=0\), the singlet condition requires perfect anti-correlation:
\begin{equation}
E_q(0)=\cos(q\pi)=(-1)^q=-1.
\end{equation}
Hence \(q\) must be odd.

\textbf{Step 2: reduction to a Chebyshev equation.}
For odd \(q\),
\begin{equation}
\cos\bigl(q(\pi-\theta)\bigr)=-\cos(q\theta).
\end{equation}
Thus exact agreement with the singlet kernel for all \(\theta\in[0,\pi]\) is equivalent to
\begin{equation}
\cos(q\theta)=\cos\theta
\qquad\hbox{for all }\theta\in[0,\pi].
\label{eq:chebyshevcondition}
\end{equation}
Since cosine is even in \(q\), set \(m=\lvert q\rvert\).  Equation~\eqref{eq:chebyshevcondition} becomes
\begin{equation}
\cos(m\theta)=\cos\theta
\qquad\hbox{for all }\theta\in[0,\pi].
\end{equation}

\textbf{Step 3: uniqueness via Chebyshev polynomials.}
Let \(T_m\) be the Chebyshev polynomial of the first kind, defined by
\begin{equation}
T_m(\cos\theta)=\cos(m\theta).
\end{equation}
With \(x=\cos\theta\), the preceding identity becomes
\begin{equation}
T_m(x)-x=0
\qquad\hbox{for all }x\in[-1,1].
\end{equation}
The left-hand side is a polynomial.  If a polynomial vanishes on a non-empty open interval, it is identically zero.  For \(m=1\), \(T_1(x)=x\), so the identity holds.  For any \(m\ge3\), the leading coefficient of \(T_m(x)\) is \(2^{m-1}\), hence \(T_m(x)-x\) has degree \(m\) and nonzero leading coefficient.  It cannot be the zero polynomial.  Therefore no odd \(m\ge3\) is allowed.

Consequently \(m=1\), i.e.
\begin{equation}
q=\pm1.
\end{equation}
The sign is the orientation convention for the Hopf line bundle.  With the standard positive Berry curvature convention for spin \(1/2\), this is written as \(q=+1\).
\end{proof}


\begin{corollary}[No topological modification of the Born rule]
\label{cor:nobornrulefromtopology}
The rigidity theorem states that exact agreement with the standard singlet kernel selects the unit Hopf charge within the Hopf-bundle lune construction.  It does not replace the Hilbert-space derivation of the Born probabilities and does not derive the Born rule from topology alone.
\end{corollary}

\begin{remark}[Why this is not a hidden-variable model]
The loop \(\Gamma_{ab}\) is not a pre-existing local variable carried by the particles.  It is a geometric representation of the joint measurement context \((\bm a,\bm b)\) in the projective Hilbert space.  The construction therefore does not restore Bell-local factorization.
\end{remark}

\section{Gauge Invariance and Abelianized \(\SU(2)\)}
\label{sec:gauge}

The connection \(A\) changes under a local section transformation \(\ket{\bm n}\mapsto e^{i\alpha(\bm n)}\ket{\bm n}\):
\begin{equation}
A\mapsto A+d\alpha.
\end{equation}
For a closed loop,
\begin{equation}
\oint_\Gamma(A+d\alpha)=\oint_\Gamma A+2\pi k,
\end{equation}
so the holonomy \(e^{i\oint A}\) is gauge invariant.  The physically meaningful quantity in Eq.~\eqref{eq:holonomycorrelation} is therefore independent of the local spinor gauge.

The word ``Abelianized'' refers to the following standard reduction.  The spin observables are non-Abelian \(\SU(2)\) objects: their algebra is
\begin{equation}
[\sigma_i,\sigma_j]=2i\epsilon_{ijk}\sigma_k.
\end{equation}
However, the geometric phase associated with a ray in \(\CP^1\) is a \(\UU(1)\) connection on the Hopf fiber.  The non-Abelian spin structure supplies the Bloch geometry and non-commuting observables; the closed-loop phase itself lies in the Abelian fiber.

Thus the construction may be summarized as
\begin{align}
\SU(2)\ \hbox{spin kinematics}
&\longrightarrow S^1\to S^3\to S^2 \nonumber\\
&\longrightarrow \UU(1)\ \hbox{closed holonomy}.
\end{align}
\section{Noether Symmetry and Conservation Law}
\label{sec:noether}

The singlet state is invariant under the diagonal action of \(\SU(2)\):
\begin{equation}
(U\otimes U)\ket{\Psi^-}=\ket{\Psi^-},
\qquad U\in\SU(2).
\label{eq:singletinvariance}
\end{equation}
Equivalently, the total spin generator annihilates the singlet,
\begin{equation}
(\bm J_A+\bm J_B)\ket{\Psi^-}=0.
\label{eq:totalspinzero}
\end{equation}
This is the relevant Noether statement: global rotational symmetry implies conservation of total angular momentum.  In correlation language,
\begin{equation}
E(R\bm a,R\bm b)=E(\bm a,\bm b),
\qquad R\in\SO(3).
\end{equation}
The holonomy representation respects the same symmetry because it depends only on the invariant relative angle \(\theta\).

\begin{remark}[Conservation without signaling]
The conservation law constrains the joint state and joint statistics.  It does not provide a controllable channel from one wing of a Bell experiment to the other.  The reduced density matrix on either side is \(\rho_A=\rho_B=\Id/2\), independent of the distant measurement setting.
\end{remark}

\section{Bell Nonlocality, No-Signaling, and the Tsirelson Bound}
\label{sec:bell}

The singlet correlations violate Bell inequalities because the joint probability distribution cannot be factorized into Bell-local hidden variables.  The holonomy reformulation does not change this fact; it repackages the correlation kernel geometrically.

The no-signaling property follows immediately from the reduced state:
\begin{equation}
\rho_A=\Tr_B(\rho_{\Psi^-})=\frac12\Id,
\qquad
\rho_B=\Tr_A(\rho_{\Psi^-})=\frac12\Id.
\end{equation}
For any local measurement \(A_{\bm a}\),
\begin{equation}
\langle A_{\bm a}\rangle=\Tr(\rho_A A_{\bm a})=0,
\end{equation}
independent of \(\bm b\).  Hence the nonlocal correlation is not a superluminal signal.

For CHSH, define
\begin{equation}
\mathcal B=A\otimes(B+B')+A'\otimes(B-B'),
\end{equation}
where \(A,A',B,B'\) are \(\pm1\)-valued spin observables.  A standard calculation gives
\begin{equation}
\mathcal B^2=4\Id-[A,A']\otimes[B,B'].
\end{equation}
Since \(\lVert[A,A']\rVert\le 2\) and \(\lVert[B,B']\rVert\le 2\),
\begin{equation}
\lVert\mathcal B\rVert\le 2\sqrt2.
\end{equation}
This is the Tsirelson bound.

A maximally violating coplanar choice is
\begin{equation}
\begin{gathered}
\bm a=(1,0,0),\qquad \bm a'=(0,1,0),\\
\bm b=\frac{1}{\sqrt2}(1,1,0),\qquad
\bm b'=\frac{1}{\sqrt2}(1,-1,0).
\end{gathered}
\end{equation}
For the singlet, the CHSH combination has magnitude \(2\sqrt2\).  In the holonomy language, each term is the real part of a closed phase; the quantum violation arises because the four phases are not values of a single joint classical probability space.  For the above maximizing settings, the corresponding Hopf phases may be written
\begin{equation}
\gamma_{ab}=\gamma_{ab'}=\gamma_{a'b}=\frac{3\pi}{4},
\qquad
\gamma_{a'b'}=\frac{\pi}{4},
\end{equation}
so that
\begin{equation}
\operatorname{Re}(U_{ab}+U_{ab'}+U_{a'b}-U_{a'b'})=-2\sqrt2.
\end{equation}
This phase display is a representation of the saturating configuration; the bound itself remains the standard operator-norm result above.

\section{Relation to the Fubini--Study Metric}
\label{sec:fs}

The projective Hilbert space \(\CP^1\simeq S^2\) carries the Fubini--Study metric.  For normalized pure states \(\ket{\psi},\ket{\phi}\), the Fubini--Study distance is
\begin{equation}
d_{\rm FS}(\psi,\phi)=\arccos |\braket{\psi}{\phi}|.
\end{equation}
For spin directions separated by Bloch angle \(\theta\),
\begin{equation}
|\braket{\bm a}{\bm b}|^2=\frac{1+\cos\theta}{2},
\qquad
 d_{\rm FS}=\frac{\theta}{2}.
\end{equation}
Thus the Bell correlation can also be written as
\begin{equation}
E(\bm a,\bm b)=-\cos(2d_{\rm FS}).
\end{equation}
Thus the same relative geometry appears as a Hilbert-space expectation value, a Fubini--Study distance, and a Hopf-bundle holonomy.

\section{Interpretation and Limits of the Construction}
\label{sec:interpretation}

The useful interpretation is narrow:
\begin{enumerate}
\item The correlation \(-\cos\theta\) is already present in standard quantum mechanics.
\item The normalized spinor space is \(S^3\), and the physical projective space is \(S^2\).
\item The Hopf connection assigns a gauge-invariant closed-loop phase to areas on \(S^2\).
\item A loop with solid angle \(2(\pi-\theta)\) gives a holonomy whose real part is \(-\cos\theta\).
\item Within this construction, exact agreement with the singlet kernel selects the unit Hopf Chern charge up to orientation.
\item The nonlocality is the usual Bell non-factorizability of the entangled state, not a new causal mechanism.
\end{enumerate}

This makes the construction a geometric representation theorem rather than a replacement theory.  It is useful because it separates what is algebraic from what is topological/geometric: the algebraic source is the singlet density matrix and Pauli algebra; the geometric carrier is the Hopf phase over the Bloch sphere.


\section{Explicit Hopf Map and Monopole Curvature}
\label{sec:hopfmap}

For a normalized spinor \(\psi=(z_1,z_2)^T\in S^3\), the Hopf map to the Bloch sphere is
\begin{equation}
\bm n(\psi)=\psi^\dagger\Pauli\psi,
\qquad
\bm n\cdot\bm n=1.
\end{equation}
The phase transformation \(\psi\mapsto e^{i\alpha}\psi\) leaves \(\bm n\) invariant, so the fiber over each Bloch vector is \(S^1\).  The Berry curvature in Eq.~\eqref{eq:berryF} is the field strength of a Dirac-monopole connection on the Bloch sphere with unit Chern number,
\begin{equation}
\frac{1}{2\pi}\int_{S^2}F=1.
\end{equation}
This is the topological origin of the sign reversal of spinors under a \(2\pi\) rotation: a closed path enclosing solid angle \(2\pi\) gives a phase \(\pi\), hence a sign change.

For the Bell kernel, the chosen loop area \(2(\pi-\theta)\) interpolates between the two limiting cases:
\begin{align}
\theta=0 &: \quad U=-1,\\
\theta=\pi &: \quad U=+1.
\end{align}
Thus the real part gives the expected limits: perfect anti-correlation at equal settings and perfect correlation at opposite settings.

\section{Joint Probabilities and Operational Content}
\label{sec:probabilities}

The holonomy expression reproduces the correlation but the operational probabilities remain the Born probabilities of the singlet state.  If \(\alpha,\beta\in\{+1,-1\}\) denote the two outcomes, then
\begin{equation}
P(\alpha,\beta|\bm a,\bm b)
=\frac14\left(1-\alpha\beta\,\bm a\cdot\bm b\right).
\label{eq:jointprob}
\end{equation}
These probabilities obey
\begin{equation}
\sum_{\beta}P(\alpha,\beta|\bm a,\bm b)=\frac12,
\qquad
\sum_{\alpha}P(\alpha,\beta|\bm a,\bm b)=\frac12,
\end{equation}
so no marginal distribution depends on the remote setting.  The correlation is recovered as
\begin{equation}
\sum_{\alpha,\beta=\pm1}\alpha\beta\,
P(\alpha,\beta|\bm a,\bm b)=-\bm a\cdot\bm b.
\end{equation}

In this sense the holonomy is not an extra observable.  It is a compact geometric encoding of the same angular dependence that appears in the joint probabilities.


\subsection{Bell-local factorization remains impossible}
\label{subsec:nolocalfactorization}

The holonomy rewriting does not weaken Bell's theorem.  Rather, it reproduces the same non-factorizable correlation kernel.

\begin{theorem}[CHSH obstruction]
For the correlation $E(a,b)=-a\cdot b$, choose coplanar settings
\begin{align}
 a&=(1,0), & a'&=(0,1),\\
 b&=\frac{1}{\sqrt2}(1,1), &
 b'&=\frac{1}{\sqrt2}(1,-1).
\end{align}
Then the CHSH combination
\begin{equation}
 S=E(a,b)+E(a,b')+E(a',b)-E(a',b')
\end{equation}
satisfies
\begin{equation}
   S=-2\sqrt2,
   \qquad |S|=2\sqrt2>2 .
\end{equation}
Therefore no Bell-local hidden-variable factorization
\begin{equation}
\begin{split}
   E(a,b)&=\int d\lambda\,\rho(\lambda)
   A(a,\lambda)B(b,\lambda),\\
   A,B&\in\{-1,+1\}.
\end{split}
\end{equation}
can reproduce the holonomy kernel.
\end{theorem}

\begin{proof}
The four correlations are
\begin{align}
E(a,b)&=E(a,b')=E(a',b)=-\frac{1}{\sqrt2},\\
E(a',b')&=+\frac{1}{\sqrt2}.
\end{align}
Hence $S=-4/\sqrt2=-2\sqrt2$. Bell-local models obey $|S|\leq2$, while the singlet/Hopf kernel attains the quantum value $2\sqrt2$.
\end{proof}

This is the precise sense in which the construction geometrizes, but does not domesticate, Bell nonlocality: the obstruction is carried by the projective Hilbert-space geometry and by the noncommuting $\SU(2)$ measurement algebra, not by a superluminal signal.

\section{Operational Recipe}
\label{sec:recipe}

Given two measurement settings \(\bm a,\bm b\):
\begin{enumerate}
\item compute the separation \(\theta\) from \(\cos\theta=\bm a\cdot\bm b\);
\item choose a closed loop \(\Gamma_{ab}\) enclosing solid angle \(2(\pi-\theta)\);
\item compute \(U(\Gamma_{ab})=\exp(i\Omega/2)\);
\item take \(E=\operatorname{Re}U\).
\end{enumerate}
The result is independent of the local spinor phase and of the loop representative, provided the oriented solid angle is fixed modulo \(4\pi\).  This is precisely the sense in which the construction is gauge invariant.

\section{Conclusion}
\label{sec:conclusion}

We have given a conservative geometric reformulation of the standard statement that Bell correlations may be expressed as Abelianized \(\SU(2)\) holonomy.  For the spin singlet,
\begin{equation}
E(\bm a,\bm b)=-\bm a\cdot\bm b
=\operatorname{Re}\exp\left(i\oint_{\Gamma_{ab}}A\right),
\end{equation}
where \(A=-i\langle\psi|d\psi\rangle\) is the Hopf/Berry connection and \(\Gamma_{ab}\) is a closed loop enclosing solid angle \(2(\pi-\theta)\).  The result is exactly equivalent to standard quantum mechanics, gauge invariant, compatible with no-signaling, and constrained by the usual Tsirelson bound.  The role of topology is not to add hidden variables, but to expose the compact spinor geometry behind the observed correlation law.  The added rigidity statement is correspondingly limited: it shows that the Hopf-bundle lune construction is fixed to unit Chern charge by exact agreement with the singlet kernel, not that topology alone replaces the Born rule or rules out other geometric representations.

\section*{Data Availability}
No datasets were generated or analyzed in this work. All calculations used in the manuscript are analytic and are displayed in the text and appendixes.

\section*{Acknowledgments}
The authors report no external funding.

\section*{Author Contributions}
A.P. and B.B. contributed to the conceptual development, analytic calculations, manuscript writing, and final verification.

\appendix

\section{Pauli Identities Used in the Text}
\label{app:pauli}

The Pauli matrices satisfy
\begin{align}
\sigma_i\sigma_j&=\delta_{ij}\Id+i\epsilon_{ijk}\sigma_k,\\
\Tr(\sigma_i)&=0,\qquad
\Tr(\sigma_i\sigma_j)=2\delta_{ij}.
\end{align}
For unit vectors \(\bm a,\bm b\),
\begin{equation}
(\bm a\cdot\Pauli)(\bm b\cdot\Pauli)
=(\bm a\cdot\bm b)\Id+i(\bm a\times\bm b)\cdot\Pauli.
\end{equation}
The eigenvalues of \(\bm a\cdot\Pauli\) are \(\pm1\).

\section{Explicit Singlet Density Matrix}
\label{app:density}

In the computational basis \((\ket{00},\ket{01},\ket{10},\ket{11})\),
\begin{equation}
\ket{\Psi^-}=\frac{1}{\sqrt2}(0,1,-1,0)^T,
\end{equation}
so
\begin{equation}
\rho_{\Psi^-}=\frac12
\begin{pmatrix}
0&0&0&0\\
0&1&-1&0\\
0&-1&1&0\\
0&0&0&0
\end{pmatrix}.
\end{equation}
Direct substitution of the Pauli matrices gives Eq.~\eqref{eq:singletdensity}.

\section{Solid Angle Convention}
\label{app:solidangle}

The solid angle in Eq.~\eqref{eq:solidangle} is defined modulo \(4\pi\), so the holonomy is defined modulo \(2\pi\).  Different loop representatives with the same oriented solid angle give the same \(\UU(1)\) holonomy.  The Bell-lune representative is therefore a gauge- and representative-independent encoding of the scalar kernel, not an additional physical trajectory.


\section{Independent Numerical Checks}
\label{app:checks}

The following values reproduce the exact identities used in the text.  For representative angles the holonomy branch is $U(\theta)=\exp[i(\pi-\theta)]$ and the correlation is $E(\theta)=\operatorname{Re}U(\theta)$:

Direct numerical checks of $E(\theta)=\operatorname{Re}e^{i(\pi-\theta)}=-\cos\theta$ are:
\begin{center}
\begin{tabular}{ccc}
\toprule
$\theta$ & $\operatorname{Re}U(\theta)$ & $-\cos\theta$ \\
\midrule
$0$ & $-1.0000000000000$ & $-1.0000000000000$ \\
$\pi/4$ & $-0.7071067811865$ & $-0.7071067811865$ \\
$\pi/2$ & $0.0000000000000$ & $0.0000000000000$ \\
$3\pi/4$ & $0.7071067811865$ & $0.7071067811865$ \\
$\pi$ & $1.0000000000000$ & $1.0000000000000$ \\
\bottomrule
\end{tabular}
\end{center}

For the CHSH settings of Sec.~\ref{subsec:nolocalfactorization}, the exact value is
\begin{align}
S&=-2.8284271247462,\\
|S|&=2.8284271247462,\\
2\sqrt2&=2.8284271247462.
\end{align}
For the Berry curvature $F=(1/2)\sin\vartheta\,d\vartheta\wedge d\phi$ on the Bloch sphere,
\begin{equation}
\begin{split}
\frac{1}{2\pi}\int_{S^2}F
&=\frac{1}{2\pi}\frac12\int_0^{2\pi}\!d\phi
  \int_0^\pi\!\sin\vartheta\,d\vartheta\\
&=1.0000000000000.
\end{split}
\end{equation}
Thus the Hopf line bundle has first Chern number $c_1=1$, as used in the text.

\section{Supported Claims and Limitations}
\label{app:limitations}

The strongest claims supported by exact calculation are:
\begin{enumerate}
\item the singlet correlation is exactly the real part of a gauge-invariant Hopf-bundle holonomy;
\item within the Hopf-bundle lune construction, the integer Chern charge is rigidly fixed to $q=\pm1$ by exact agreement with the singlet kernel, with the sign set by orientation convention;
\item the construction preserves no-signaling, the Born rule, and the Tsirelson bound;
\item the construction cannot be a Bell-local hidden-variable model.
\end{enumerate}
The following claims are not supported without adding new physical postulates: new experimental predictions, a violation or replacement of standard quantum mechanics, a causal superluminal mechanism, a derivation of the Born rule from topology alone, or uniqueness over all possible geometric representations of the singlet kernel.

\begin{thebibliography}{99}

\bibitem{Bell1964}
J.~S. Bell,
On the Einstein Podolsky Rosen paradox,
Physics \textbf{1}, 195 (1964).

\bibitem{CHSH1969}
J.~F. Clauser, M.~A. Horne, A.~Shimony, and R.~A. Holt,
Proposed experiment to test local hidden-variable theories,
Phys. Rev. Lett. \textbf{23}, 880 (1969).

\bibitem{Tsirelson1980}
B.~S. Tsirelson,
Quantum generalizations of Bell's inequality,
Lett. Math. Phys. \textbf{4}, 93 (1980).

\bibitem{Berry1984}
M.~V. Berry,
Quantal phase factors accompanying adiabatic changes,
Proc. R. Soc. Lond. A \textbf{392}, 45 (1984).

\bibitem{Simon1983}
B.~Simon,
Holonomy, the quantum adiabatic theorem, and Berry's phase,
Phys. Rev. Lett. \textbf{51}, 2167 (1983).

\bibitem{Nakahara2003}
M.~Nakahara,
\textit{Geometry, Topology and Physics}, 2nd ed.
(Taylor \& Francis, Boca Raton, 2003).

\bibitem{NielsenChuang}
M.~A. Nielsen and I.~L. Chuang,
\textit{Quantum Computation and Quantum Information}
(Cambridge University Press, Cambridge, 2000).

\bibitem{BengtssonZyczkowski}
I.~Bengtsson and K.~{\.{Z}}yczkowski,
\textit{Geometry of Quantum States}, 2nd ed.
(Cambridge University Press, Cambridge, 2017).

\end{thebibliography}

\end{document}
